\part{Categorical Interpretations}
\label{part:X}

Part X interprets economic processes as morphisms and states what the interpretation preserves. Its stated payoff is that parallel composition fails to be independent under shared sinks, and the structure needed to repair that is identified.

\begin{housing}Repair processes become morphisms between building interfaces, and accounting becomes a functor that forgets the backlog.\end{housing}

\chapter{Economic Processes as Morphisms}\label{ch:46}
Takes specified economic interfaces or states as objects and admissible transformations, carrying explicit resource and obligation effects, as morphisms. Identity and composition are defined before any categorical claim. The construction begins with the syntactic congruence of the admissible-history language, and where ordinary composition is inadequate the chapter states the extra structure required.
\chspec{\Pref{P01}, \Pref{P05}, \Pref{P18}.}
\scaffold

\chapter{Parallel Production and Monoidal Structure}\label{ch:47}
Represents processes performed alongside one another through a monoidal structure. Shared labour, ecological sinks, or scarce machinery complicate the interpretation of independence. String-diagram notation stays optional, and algebraic presentations suffice for the arguments.
\chspec{\Pref{P09}, \Pref{P10}, \Pref{P18}.}
\scaffold

\chapter{Resources, Obligations, and Enriched Categories}\label{ch:48}
Motivates enrichment over appropriate ordered structures by costs, delays, risks, and resource requirements. Develops multidimensional cost composition and identifies when a single numerical distance suppresses economically relevant distinctions.
\chspec{\Pref{P13}, \Pref{P24}.}
\scaffold

\chapter{Accounting as a Structure-Preserving Map}\label{ch:49}
Connects operational processes to accounting representations by functors. The central question is what information an accounting map preserves, identifies, or loses. Two economically different processes can have the same financial image, which gives a precise form of accounting insufficiency relative to a stated task.
\chspec{\Pref{P11}, \Pref{P12}, \Pref{P19}, \Pref{P29}.}
\scaffold

\chapter{Open Systems, Institutions, and Compositional Guarantees}\label{ch:50}
Describes households, firms, and infrastructures through interfaces in an open-systems formalism. Develops an assume-guarantee treatment of institutional composition and identifies conditions under which local guarantees support a global continuation claim. Categorical language organizes the proof obligations while empirical premises remain independently necessary.
\chspec{\Pref{P04}, \Pref{P23}, \Pref{P27}.}
\scaffold
